\begin{problem}[Naive constant versus locally constant]
For an abelian group $A$, compare the naive constant presheaf $\underline A^{pre}$ with the presheaf $\mathcal L_A$ of locally constant $A$-valued functions.  Construct a natural morphism $\underline A^{pre}\to\mathcal L_A$ and determine when it is an isomorphism on an open $U$.
\end{problem}

\paragraph{Why this problem matters.}
This comparison isolates the connected-component information the naive constant presheaf forgets.

\paragraph{Useful ideas before starting.}
Send $a\in A$ to the constant function with value $a$.

\paragraph{Guided hints.}
Send $a\in A$ to the constant function with value $a$.

\begin{solution}
The maps $A\to\mathcal L_A(U)$ sending $a$ to the constant function are compatible with restriction.  On nonempty $U$ they are injective.  They are surjective exactly when every locally constant function on $U$ is constant, in particular when $U$ is connected.  On a disconnected open, different values on components give sections outside the image.
\end{solution}

\paragraph{What happened mathematically.}
Locally constant functions permit independent constants on components.

\paragraph{Extensions.}
Study what happens on locally connected spaces.

\paragraph{Provenance.}
Legacy-derived from the AG1 constant-presheaf example.
